\documentclass[11pt]{article}

\usepackage[margin=1in]{geometry}
\usepackage{amsmath,amssymb,amsthm}
\usepackage{bm}
\usepackage{booktabs}
\usepackage{hyperref}

\newcommand{\vect}[1]{\bm{#1}}
\newcommand{\tens}[1]{\mathsf{#1}}
\newcommand{\ez}{\hat{\vect{e}}_z}
\newcommand{\Rm}{\tens{R}}
\newcommand{\Gl}{\vect{G}}
\newcommand{\Gb}{\vect{G}_{b}}
\newcommand{\T}{^{\mathsf{T}}}

\theoremstyle{remark}
\newtheorem*{remark}{Remark}

\title{Physics of \texttt{sim\_symmetric.jsx}:\\
Rigid-Body Dynamics in a Harmonic Thermophoretic Trap}
\author{Chin Lab, University of Chicago}
\date{\today}

\begin{document}
\maketitle

\begin{abstract}
This note sets out the model solved by the interactive simulator
\texttt{Theory/sim\_symmetric.jsx}. A rigid body of mass $M$ and inertia
tensor $\tens{I}$ moves through a quiescent gas under gravity. It feels a
thermophoretic force and torque that are linear in the local
log-temperature gradient $\nabla\ln T$, plus linear translational and
rotational drag. The temperature field is a phenomenological harmonic
trap. Each of the four linear-response tensors is a constant symmetric
$3\times3$ matrix in the body frame. The 13-dimensional state (position,
velocity, unit quaternion, body angular velocity) is integrated with
fixed-step fourth-order Runge--Kutta. We give the equations exactly as
implemented, derive the equilibrium and the small-oscillation behaviour,
and list the modelling assumptions and their consequences.
\end{abstract}

%=======================================================================
\section{Frames and state variables}
%=======================================================================

We use two frames: the \emph{lab frame}
$\{\hat{\vect{e}}_x,\hat{\vect{e}}_y,\hat{\vect{e}}_z\}$, with $\ez$
pointing up (against gravity), and the \emph{body frame}, which is attached
to the particle and whose axes are its principal axes of inertia. The
rotation matrix $\Rm\in SO(3)$ maps body-frame components to lab-frame
components:
\begin{equation}
  \vect{a}_{\text{lab}} = \Rm\,\vect{a}_b, \qquad
  \vect{a}_b = \Rm\T \vect{a}_{\text{lab}} .
\end{equation}

The state vector $\vect{s}\in\mathbb{R}^{13}$ is
\begin{equation}
  \vect{s} = \bigl(\,\underbrace{\vect{x}}_{3},\;
                  \underbrace{\vect{v}}_{3},\;
                  \underbrace{q}_{4},\;
                  \underbrace{\vect{\omega}_b}_{3}\,\bigr),
\end{equation}
where
\begin{itemize}
  \item $\vect{x}$ is the centre-of-mass position (lab frame),
  \item $\vect{v}=\dot{\vect{x}}$ is the centre-of-mass velocity (lab frame),
  \item $q=(q_0,\vect{q})$, with $\vect{q}=(q_1,q_2,q_3)$, is a unit
        quaternion ($|q|=1$) that encodes the orientation, and
  \item $\vect{\omega}_b$ is the angular velocity, expressed in the body frame.
\end{itemize}

\paragraph{Rotation matrix.} The quaternion determines the rotation
matrix through
\begin{equation}
  \Rm(q) = \bigl(q_0^2-|\vect{q}|^2\bigr)\tens{1}
         + 2\,\vect{q}\vect{q}\T + 2q_0[\vect{q}]_\times ,
  \qquad
  [\vect{q}]_\times =
  \begin{pmatrix}
    0 & -q_3 & q_2\\ q_3 & 0 & -q_1\\ -q_2 & q_1 & 0
  \end{pmatrix}.
\end{equation}
Written out using $|q|=1$, this is
\begin{equation}
  \Rm =
  \begin{pmatrix}
    1-2(q_2^2+q_3^2) & 2(q_1q_2-q_0q_3) & 2(q_1q_3+q_0q_2)\\
    2(q_1q_2+q_0q_3) & 1-2(q_1^2+q_3^2) & 2(q_2q_3-q_0q_1)\\
    2(q_1q_3-q_0q_2) & 2(q_2q_3+q_0q_1) & 1-2(q_1^2+q_2^2)
  \end{pmatrix}.
\end{equation}

\paragraph{Quaternion kinematics.} The orientation evolves as
\begin{equation}
  \dot q = \tfrac12\, q\otimes(0,\vect{\omega}_b)
  \quad\Longleftrightarrow\quad
  \dot q_0 = -\tfrac12\,\vect{q}\cdot\vect{\omega}_b, \qquad
  \dot{\vect{q}} = \tfrac12\bigl(q_0\vect{\omega}_b + \vect{q}\times\vect{\omega}_b\bigr).
  \label{eq:qdot}
\end{equation}
The body-frame angular velocity multiplies on the right in
Eq.~\eqref{eq:qdot}. This is what makes $\Rm$ a body-to-lab map.

\paragraph{Tilt.} The simulator reports the \emph{tilt}, defined as the angle
between the body symmetry axis $\hat{\vect{e}}_3^{\,b}$ and the vertical:
\begin{equation}
  \theta = \arccos\bigl(\ez\cdot\Rm\,\hat{\vect{e}}_3\bigr) = \arccos R_{33}.
\end{equation}

%=======================================================================
\section{Temperature field}
%=======================================================================

The gas temperature enters only through the dimensionless-per-length
vector field
\begin{equation}
  \Gl(\vect{x}) \equiv \nabla\ln T = \frac{\nabla T}{T}.
\end{equation}
The simulator takes this field to be affine in position. In the lab frame,
\begin{equation}
  \boxed{\;
  \Gl(\vect{x}) = \tens{A}\,\vect{x} + \Gl_0, \qquad
  \tens{A} = \operatorname{diag}(\alpha_r,\alpha_r,\alpha_z), \qquad
  \Gl_0 = -\frac{Mg}{\gamma_\ast}\,\ez \;}
  \label{eq:field}
\end{equation}
Component by component,
\begin{equation}
  \Gl(\vect{x}) = \bigl(\alpha_r x,\ \alpha_r y,\ \alpha_z z - Mg/\gamma_\ast\bigr).
\end{equation}
The corresponding log-temperature profile is a paraboloid with a linear
tilt:
\begin{equation}
  \ln\frac{T(\vect{x})}{T_0}
  = \tfrac12\alpha_r\,(x^2+y^2) + \tfrac12\alpha_z z^2 - \frac{Mg}{\gamma_\ast}\,z .
\end{equation}

\begin{itemize}
  \item \textbf{Uniform part $\Gl_0$.} $\Gl_0$ points down, so the gas is
        hotter below and colder above, as with a hot bottom plate. Its
        magnitude is fixed by the particle itself: the code sets
        $|\Gl_0| = Mg/\gamma_\ast$ so that, for an isotropic response, the
        thermophoretic force at the origin exactly cancels gravity
        (Sec.~\ref{sec:equilibrium}). The trap is therefore always
        centred at $\vect{x}=0$. Changing $M$ or $\Gamma$ does not move the
        equilibrium; it only rescales the background gradient.
  \item \textbf{Curvature $\tens{A}$.} The constants $\alpha_r$ and
        $\alpha_z$ (units m$^{-2}$) set the radial and axial trap
        stiffness. They are free phenomenological parameters.
  \item \textbf{Normalising constant $\gamma_\ast$.} The code uses
        $\gamma_\ast = \Gamma_{22}$, the $(2,2)$ entry of the
        thermophoretic tensor (\texttt{Gamma[3]} in the packed
        storage of Sec.~\ref{sec:tensors}). See caveat~\ref{rem:gamma22} in Sec.~\ref{sec:caveats}.
\end{itemize}

The gas is at rest everywhere, and the temperature field does not depend
on time.

%=======================================================================
\section{Linear-response tensors}
\label{sec:tensors}
%=======================================================================

Four constant $3\times3$ tensors, defined in the body frame, set how the
body couples to the gas. Let $\vect{v}_b=\Rm\T\vect{v}$ be the body-frame
velocity and $\Gb=\Rm\T\Gl(\vect{x})$ the gradient seen by the body. Then
\begin{align}
  \vect{F}_b^{\text{drag}} &= -\,\tens{\Xi}\,\vect{v}_b
      && \text{translational drag,}\\
  \vect{F}_b^{\text{T}}    &= -\,\tens{\Gamma}\,\Gb
      && \text{thermophoretic force,}\\
  \vect{\tau}_b^{\text{T}} &= -\,\tens{\Lambda}\,\Gb
      && \text{thermophoretic torque,}\\
  \vect{\tau}_b^{\text{drag}} &= -\,\tens{C}_r\,\vect{\omega}_b
      && \text{rotational drag.}
\end{align}
Take $\tens{\Gamma}=\gamma\tens{1}$ with $\gamma>0$ as an example. The force
$-\gamma\nabla\ln T$ then points from hot to cold, as expected for
thermophoresis.

\paragraph{Symmetry by construction.} The code stores each tensor as the
six independent entries of a symmetric matrix,
\begin{equation}
  \tens{t} \;\longleftrightarrow\; (t_{11},t_{12},t_{13},t_{22},t_{23},t_{33}),
  \qquad
  \tens{t} =
  \begin{pmatrix}
    t_{11} & t_{12} & t_{13}\\
    t_{12} & t_{22} & t_{23}\\
    t_{13} & t_{23} & t_{33}
  \end{pmatrix}
  = \tens{t}\T .
\end{equation}
This gives $4\times6=24$ independent response parameters. The rationale
given in the code is:
\begin{itemize}
  \item $\tens{\Xi}$ and $\tens{C}_r$ are resistance tensors. For
        low-Reynolds-number flow, the Lorentz reciprocal theorem (the
        hydrodynamic form of Onsager reciprocity) makes them symmetric.
        The second law adds positive (semi)definiteness, because the
        dissipated power is
        \begin{equation}
          P_{\text{diss}} = \vect{v}_b\cdot\tens{\Xi}\,\vect{v}_b
                          + \vect{\omega}_b\cdot\tens{C}_r\,\vect{\omega}_b \;\ge 0 .
        \end{equation}
        The simulator enforces symmetry but does \emph{not} enforce
        positive-definiteness. The user must pick physically sensible
        values.
  \item $\tens{\Gamma}$ and $\tens{\Lambda}$ are taken to be symmetric,
        motivated by Onsager reciprocity. This is a modelling choice, and
        its consequences are discussed in Sec.~\ref{sec:caveats}.
  \item $\tens{\Lambda}$ is a pseudotensor: it maps a polar vector
        ($\nabla\ln T$) to an axial vector (torque). Its symmetric part
        can be nonzero only for bodies without enough improper symmetry;
        in particular, its trace requires a chiral body. The default
        $\tens{\Lambda}=0$ is correct for a right circular cylinder.
\end{itemize}

\paragraph{Lab-frame form.} In the lab frame the tensors become
orientation dependent:
\begin{equation}
  \tens{\Gamma}'(q) \equiv \Rm\,\tens{\Gamma}\,\Rm\T, \qquad
  \tens{\Xi}'(q) \equiv \Rm\,\tens{\Xi}\,\Rm\T .
\end{equation}
An isotropic tensor $\tens{t}=t\tens{1}$ is unchanged by this rotation.

%=======================================================================
\section{Equations of motion}
%=======================================================================

\subsection{Translation}

Newton's second law for the centre of mass, with the forces evaluated in
the body frame and then rotated back to the lab frame, reads
\begin{equation}
  M\dot{\vect{v}}
  = \Rm\Bigl[-\tens{\Gamma}\,\Rm\T\Gl(\vect{x}) - \tens{\Xi}\,\Rm\T\vect{v}\Bigr]
    - Mg\,\ez .
  \label{eq:trans}
\end{equation}
Equivalently,
\begin{equation}
  M\ddot{\vect{x}}
  = -\tens{\Gamma}'\bigl(\tens{A}\vect{x}+\Gl_0\bigr)
    - \tens{\Xi}'\,\dot{\vect{x}} - Mg\,\ez .
\end{equation}

\subsection{Rotation}

Euler's equations in the principal body frame, with
$\tens{I}=\operatorname{diag}(I_1,I_2,I_3)$, are
\begin{equation}
  \tens{I}\,\dot{\vect{\omega}}_b
  = -\,\vect{\omega}_b\times\bigl(\tens{I}\,\vect{\omega}_b\bigr)
    - \tens{\Lambda}\,\Gb - \tens{C}_r\,\vect{\omega}_b ,
  \qquad \Gb = \Rm\T\Gl(\vect{x}).
  \label{eq:rot}
\end{equation}
Written out in components, the gyroscopic term gives the familiar form
\begin{equation}
  I_1\dot\omega_1 = (I_2-I_3)\,\omega_2\omega_3 + \tau_1,\quad
  I_2\dot\omega_2 = (I_3-I_1)\,\omega_3\omega_1 + \tau_2,\quad
  I_3\dot\omega_3 = (I_1-I_2)\,\omega_1\omega_2 + \tau_3,
\end{equation}
with $\vect{\tau} = -\tens{\Lambda}\Gb - \tens{C}_r\vect{\omega}_b$.
The torques are taken about the centre of mass, and the thermophoretic
force is assumed to act through that point.

\subsection{Complete system}

Together with Eq.~\eqref{eq:qdot}, the model is the closed autonomous
system
\begin{equation}
  \boxed{
  \begin{aligned}
    \dot{\vect{x}} &= \vect{v},\\
    M\dot{\vect{v}} &= -\Rm\bigl(\tens{\Gamma}\Gb + \tens{\Xi}\,\Rm\T\vect{v}\bigr) - Mg\,\ez,\\
    \dot q &= \tfrac12\,q\otimes(0,\vect{\omega}_b),\\
    \tens{I}\dot{\vect{\omega}}_b &= -\vect{\omega}_b\times\tens{I}\vect{\omega}_b
        - \tens{\Lambda}\Gb - \tens{C}_r\vect{\omega}_b,
  \end{aligned}
  \qquad \Gb = \Rm(q)\T\bigl(\tens{A}\vect{x}+\Gl_0\bigr).}
\end{equation}
The translational and rotational sectors are coupled in two ways:
\begin{enumerate}
  \item through $\Rm(q)$ in the force, if $\tens{\Gamma}$ or
        $\tens{\Xi}$ is anisotropic, and
  \item through $\Gb(\vect{x},q)$ in the torque, if $\tens{\Lambda}\neq0$.
\end{enumerate}

%=======================================================================
\section{Equilibrium and small oscillations}
\label{sec:equilibrium}
%=======================================================================

\subsection{Isotropic response}

Take $\tens{\Gamma}=\gamma\tens{1}$ and $\tens{\Xi}=\xi\tens{1}$, so that
$\gamma_\ast=\gamma$. Then $\tens{\Gamma}'=\gamma\tens{1}$ for every
orientation. The translational equation decouples from the rotation and
is \emph{exactly} linear:
\begin{equation}
  M\ddot{\vect{x}} = -\gamma\tens{A}\vect{x} - \xi\dot{\vect{x}}
  + \underbrace{\bigl(-\gamma\Gl_0 - Mg\,\ez\bigr)}_{=\,0}.
\end{equation}
Gravity is cancelled identically by the background gradient. Each axis
$i\in\{x,y,z\}$ is an independent damped harmonic oscillator:
\begin{equation}
  \ddot x_i + 2\zeta_i\omega_{0,i}\,\dot x_i + \omega_{0,i}^2 x_i = 0,
  \qquad
  \omega_{0,i}^2 = \frac{\gamma\alpha_i}{M},
  \qquad
  \zeta_i = \frac{\xi}{2\sqrt{M\gamma\alpha_i}} ,
\end{equation}
where $\alpha_{x}=\alpha_y=\alpha_r$ and $\alpha_z=\alpha_z$. The origin is
asymptotically stable if $\gamma\alpha_i>0$ and $\xi>0$. The translational
energy
\begin{equation}
  E_{\text{tr}} = \tfrac12 M|\vect{v}|^2
  + \tfrac12\gamma\bigl[\alpha_r(x^2+y^2)+\alpha_z z^2\bigr]
\end{equation}
is a Lyapunov function with $\dot E_{\text{tr}} = -\xi|\vect{v}|^2$.

\paragraph{Default parameters.} With $M=1$, $\gamma=5$,
$\alpha_r=\alpha_z=0.5$, and $\xi=2$, all three axes have
\begin{equation}
  \omega_0=\sqrt{2.5}\approx1.58~\text{s}^{-1},\quad
  \zeta\approx0.63,\quad
  \omega_d=\omega_0\sqrt{1-\zeta^2}\approx1.22~\text{s}^{-1},\quad
  \text{decay rate } \frac{\xi}{2M}=1~\text{s}^{-1}.
\end{equation}
The motion is underdamped, with period $2\pi/\omega_d\approx5.1$\,s. For
example, starting from rest,
$z(t)=z_0\,e^{-t}\bigl[\cos\omega_d t+\omega_d^{-1}\sin\omega_d t\bigr]$
with $z_0=0.2$. The simulator output matches this to its printed precision
($z=0.0815,\,-0.0067,\,0.0012$ at $t=1,2,5$\,s).

\subsection{Anisotropic response}

For a general $\tens{\Gamma}$, the equilibrium at fixed orientation $q$
satisfies
\begin{equation}
  \tens{\Gamma}'\bigl(\tens{A}\vect{x}^\ast + \Gl_0\bigr) = -Mg\,\ez
  \quad\Longrightarrow\quad
  \vect{x}^\ast = Mg\,\tens{A}^{-1}\Bigl[\frac{\ez}{\gamma_\ast}
                 - \tens{\Gamma}'^{-1}\ez\Bigr].
\end{equation}
The trap centre lies at the origin only when $\ez$ is an eigenvector of
$\tens{\Gamma}'$ with eigenvalue $\gamma_\ast$. Otherwise the equilibrium
is displaced, and the displacement depends on orientation. Near
equilibrium the effective stiffness matrix is
\begin{equation}
  \tens{K} = \tens{\Gamma}'\tens{A}.
\end{equation}
$\tens{K}$ is symmetric only when $\tens{\Gamma}'$ and $\tens{A}$ commute,
which holds automatically when $\alpha_r=\alpha_z$. If they do not
commute, the restoring force has a nonzero curl
($\nabla\times\vect{F}\neq0$). It is then a non-conservative ``curl
force'' that can drive circulating orbits, so $E_{\text{tr}}$ no longer
works as a Lyapunov function.

%=======================================================================
\section{Rotational dynamics}
%=======================================================================

With $\tens{\Lambda}=0$ there is \emph{no orientational restoring torque}:
every orientation is neutrally stable. Take isotropic rotational drag
$\tens{C}_r=c\tens{1}$ and an axisymmetric body with
$I_1=I_2\equiv I_\perp$. Euler's equations can then be solved in closed
form. The spin about the symmetry axis decays exponentially,
\begin{equation}
  \omega_3(t) = \omega_3(0)\,e^{-ct/I_3}.
\end{equation}
The transverse component $\omega_\perp\equiv\omega_1+i\omega_2$ obeys
\begin{equation}
  \dot\omega_\perp = \bigl(i\Omega(t) - c/I_\perp\bigr)\,\omega_\perp,
  \qquad
  \Omega(t) = \frac{I_3-I_\perp}{I_\perp}\,\omega_3(t),
\end{equation}
which integrates to
\begin{equation}
  \omega_\perp(t) = \omega_\perp(0)\,
  \exp\!\Bigl[-\frac{c\,t}{I_\perp}
  + i\,\frac{(I_3-I_\perp)I_3}{I_\perp c}\,\omega_3(0)\bigl(1-e^{-ct/I_3}\bigr)\Bigr].
\end{equation}
So $\vect{\omega}_b$ performs a damped body-frame precession and falls to
zero. The rotational energy $\tfrac12\vect{\omega}_b\cdot\tens{I}\vect{\omega}_b$
decreases at the rate $c|\vect{\omega}_b|^2$. The orientation then
\emph{freezes} at whatever value it has reached.

\paragraph{Default parameters.} With $\tens{I}=\operatorname{diag}(0.02,0.02,0.008)$
and $c=0.5$, the spin-down times are $I_\perp/c=0.04$\,s and $I_3/c=0.016$\,s.
Both are much shorter than the translational time scale. The total
transverse rotation angle is only about $|\omega_\perp(0)|\,I_\perp/c\approx0.02$\,rad,
so the tilt stays at its initial $30^\circ$ to the displayed precision.

\paragraph{Nonzero $\tens{\Lambda}$.} At the trap centre, a body with
$\tens{\Lambda}\neq0$ still feels the torque
$-\tens{\Lambda}\Rm\T\Gl_0 = (Mg/\gamma_\ast)\,\tens{\Lambda}\,\Rm\T\ez$,
because $\Gl_0\neq0$. The body therefore generically does not come to
rest. A chiral body, for example, can be driven into steady rotation, with
the thermophoretic torque balanced by rotational drag.

%=======================================================================
\section{Numerical method}
%=======================================================================

\begin{itemize}
  \item \textbf{Integrator.} Classical fixed-step RK4 on
        $\dot{\vect{s}}=\vect{f}(\vect{s})$, with
        $\Delta t = t_{\text{end}}/N$. The default is
        $t_{\text{end}}=20$\,s and $N=1200$, which gives $\Delta t=1/60$\,s.
  \item \textbf{Quaternion normalisation.} $q$ is renormalised inside every
        right-hand-side evaluation and again after every step. This removes
        the drift off $|q|=1$ that the unconstrained RK4 update would
        otherwise accumulate.
  \item \textbf{Initial conditions.} The initial orientation is a rotation
        by $\theta_0$ about $\hat{\vect{e}}_x$:
        $q(0)=(\cos\tfrac{\theta_0}{2},\sin\tfrac{\theta_0}{2},0,0)$.
        The user also sets $\vect{x}_0$, $\vect{v}_0$, and $\vect{\omega}_{b,0}$.
  \item \textbf{Stability.} Explicit RK4 needs $|\lambda|\Delta t\lesssim2.8$
        for real, negative eigenvalues $\lambda$. The fastest default mode
        is the axial spin-down, with $c/I_3=62.5$\,s$^{-1}$, giving
        $|\lambda|\Delta t\approx1.04$. This is stable, though the fast
        rotational transient is resolved by only a few steps. Making
        $I_3$ smaller or $c$ larger by a factor of about~3 would push the
        integration past the stability limit.
\end{itemize}

%=======================================================================
\section{Modelling assumptions and caveats}
\label{sec:caveats}
%=======================================================================

\begin{enumerate}
  \item \textbf{Linear response, constant coefficients.} All four tensors
        are independent of position, temperature, pressure, and velocity.
        Real thermophoretic coefficients depend on the local gas state
        (for example, $T$, $p$, and the Knudsen number) and vary across
        the cell.

  \item \textbf{Missing cross-couplings.} The full $6\times6$ resistance
        matrix of a rigid body has an off-diagonal translation--rotation
        block, which is nonzero for chiral bodies and for bodies whose
        reference point is not the centre of reaction. The model omits it.
        It likewise omits any force or torque proportional to higher
        gradients of $T$.

  \item \textbf{Symmetric $\tens{\Lambda}$ excludes alignment torques.}
        Consider a polar but achiral body, such as a cone or a Janus
        particle, with symmetry group $C_{\infty v}$ and axis
        $\hat{\vect{n}}$. The most general $\tens{\Lambda}$ allowed by
        this symmetry is purely \emph{antisymmetric}, giving the aligning
        torque $\vect{\tau}\propto\hat{\vect{n}}\times\nabla\ln T$.
        Imposing $\tens{\Lambda}=\tens{\Lambda}\T$ forbids this torque. A
        symmetric $\tens{\Lambda}$ describes chirality-type couplings
        instead. For example, $D_\infty$ (helix-like) bodies give
        $\tens{\Lambda}=\operatorname{diag}(a,a,b)$.

  \item \textbf{Phenomenological trap.}
        \label{rem:gammaconstraint}
        $\alpha_r$ and $\alpha_z$ are free parameters. If $T$ instead
        obeyed source-free steady conduction with $\kappa\propto T^n$,
        then $\nabla\!\cdot(T^{n+1}\nabla\ln T)=0$, which implies
        \begin{equation}
          \nabla^2\ln T = -(n+1)\,|\nabla\ln T|^2 .
        \end{equation}
        At the trap centre this requires
        $2\alpha_r+\alpha_z=-(n+1)(Mg/\gamma)^2<0$ for $n>-1$. So a
        constant isotropic $\gamma$ acting on a purely conductive field
        cannot confine the particle along all three axes at once; this is
        an Earnshaw-type obstruction. The harmonic trap should therefore
        be read as an effective description that absorbs whatever
        provides confinement in the experiment, such as spatially varying
        coefficients, boundaries, or convection.

  \item \textbf{Choice of $\gamma_\ast$.}
        \label{rem:gamma22}
        The code sets $\gamma_\ast=\Gamma_{22}$ (\texttt{Gamma[3]}). But
        gravity is along $\ez$, and for an upright body with diagonal
        $\tens{\Gamma}$ the component that cancels gravity is
        $\Gamma_{33}$ (\texttt{Gamma[5]}). The two coincide for isotropic
        $\tens{\Gamma}$, including the defaults. For
        $\tens{\Gamma}=\operatorname{diag}(\Gamma_{11},\Gamma_{22},\Gamma_{33})$
        with $\Gamma_{22}\neq\Gamma_{33}$, the equilibrium of an upright
        body moves to
        \begin{equation}
          z^\ast=\frac{Mg}{\alpha_z}\Bigl(\frac{1}{\Gamma_{22}}-\frac{1}{\Gamma_{33}}\Bigr).
        \end{equation}

  \item \textbf{Positivity not enforced.} Nothing in the code prevents
        indefinite $\tens{\Xi}$ or $\tens{C}_r$. Such values would inject
        energy and are unphysical.

  \item \textbf{Principal axes.} $\tens{I}$ is diagonal in the same body
        frame in which the response tensors are specified. If the
        principal axes of inertia do not coincide with the frame in which
        the tensors are known, the tensors must first be rotated into the
        principal frame.
\end{enumerate}

%=======================================================================
\appendix
\section{Default parameters}
%=======================================================================

\begin{center}
\begin{tabular}{@{}lll@{}}
\toprule
Quantity & Symbol & Default (as in code) \\
\midrule
Mass & $M$ & $1$ kg \\
Gravity & $g$ & $9.81$ m\,s$^{-2}$ \\
Inertia & $\tens{I}$ & $\operatorname{diag}(0.02,0.02,0.008)$ kg\,m$^2$ \\
Translational drag & $\tens{\Xi}$ & $2\,\tens{1}$ \\
Thermophoretic force & $\tens{\Gamma}$ & $5\,\tens{1}$ \\
Thermophoretic torque & $\tens{\Lambda}$ & $0$ \\
Rotational drag & $\tens{C}_r$ & $0.5\,\tens{1}$ \\
Trap curvatures & $\alpha_r,\ \alpha_z$ & $0.5,\ 0.5$ \\
Initial position & $\vect{x}_0$ & $(0.5,\,0.3,\,0.2)$ m \\
Initial velocity & $\vect{v}_0$ & $0$ \\
Initial body angular velocity & $\vect{\omega}_{b,0}$ & $(0,\,0.5,\,2)$ rad\,s$^{-1}$ \\
Initial tilt (about $\hat{\vect{e}}_x$) & $\theta_0$ & $30^\circ$ \\
Duration, steps & $t_{\text{end}},\ N$ & $20$ s, $1200$ \\
\bottomrule
\end{tabular}
\end{center}

Note that the README lists $\tens{\Xi}=5\,\tens{1}$, but the code default
is $2\,\tens{1}$. With $\xi=5$, the translational damping ratio would be
$\zeta\approx1.58$, which is overdamped, rather than $0.63$.

\end{document}
